% Page 056

\seite{56}

\abschnitt{Schuljahr 1919–20}
\pstart
\margmark{Entlassen wu\\der Landwirt\\Knaben. Die\\32 Mädchen.}%
rden vierzig 2 Knaben, Ostern 1 Mädchen. Alle widmen sich\ob
schaft. Aufgenommen wurden 8 Kinder, 6 Mädchen, 2\ob
Schülerzahl beträgt so jetzt 44 Schüler, davon 12 Knaben,\ob

\pend
\abschnitt{Kriegsfolge 5}
\pstart

\margmark{Der Krieg is\\Helden werde\\eingetroffen\\verloren, 4\\zurück, 3 an\\Heimat anlan}%
t beendet, das Dorf hat 7 Gefallene zu beklagen. Die tapferen\ob
n wir nicht vergessen. Die anderen Krieger sind in die Heimat\ob
, d.\ h. ein Jüngling von 23 Jahren als Krüppel, ein Bein hat er\ob
aus dem Dorfe waren in Gefangenschaft, einer ist in englischer\ob
dere sind in französischer. Alle hoffen sie auf bald in der\ob
gen zu können.

\pend
\abschnitt{Ernte}
\pstart

\margmark{Die Gesamter\\war verhältn\\eingebracht.\\preise als B}%
nte lieferte stellenweise einen geringen Ertrag. Die Getreideernte\ob
ismäßig galt gegen letztes Jahr mehr Kartoffelfrüchte viel\ob
 Die Preise im allgemeinen seigen fortwährlig, sowohl Vieh\ohb
utterpreise.

\pend
\abschnitt{Schuljahr 1920–21}
\pstart

\margmark{Vorzeitig wu\\1920 1 Mädch\\wirtschaft.\\ben. Die Sch}%
rde Ostern Herbst 1919 ein Schüler entlassen, Ostern\ob
en und 3 Knaben, sämtliche widmen sich der Land\ohb
Aufgenommen wurden an Kinder: 4 Mädchen, 3 Kna\ohb
ülerzahl beträgt 24, 29 Mädchen, 15 Knaben.

\pend
\abschnitt{Ernte}
\pstart

\margmark{In der hiesi\\wohl das Get\\tet hatte. D\\auch für Kar\\stark gefall\\sperre sollt\\besetzten Ge}%
gen Gemeinde fiel allgemein befriedigend aus, so\ohb
reide, wie auch die\unsicher{} Kartoffelernte besser, als man erwar\ohb
ie Witterung war eine besonders günstige, namentlich\ob
toffeln. Die Preise sind für Getreide und Vieh sehr\ob
en, es ist schlecht, infolge der Einführung Klauenseuche-Grenz\ohb
e im Herbste große Märkte ausfallen, auch sollten von den\ob
bieten Lebensmittelzufuhr sind noch wie nur immer sehr schwach.
\pend
